\documentclass[11pt,a4paper]{article}
\usepackage[margin=2.0cm]{geometry}
\usepackage{lmodern}
\usepackage[T1]{fontenc}
\usepackage{microtype}
\usepackage{hyperref}
\usepackage{enumitem}
\usepackage{booktabs}
\usepackage{tikz}
\usepackage{listings}
\usepackage{xcolor}
\usetikzlibrary{positioning, arrows.meta, calc, shapes.geometric, fit, matrix}

\title{Experiment Runner Framework -- design + UML class diagram}
\author{2026-06-04 (Aiko / Csaba)}
\date{}

\begin{document}
\maketitle

\section*{Scope}
Replace 221 launcher scripts in \texttt{signedkan\_wip/experiments/}
with one runner framework + YAML configs. CLAUDE.md \S6.5 \#3 has
mandated this since 2026-05-11; this is the implementation.

The full design narrative is in
\texttt{runner\_architecture.md} (in the same directory).
This LaTeX wrapper carries the formal UML class diagram below.

\section*{Class diagram}
\begin{figure}[h]
  \centering
  % Experiment Runner Framework -- UML class diagram (TikZ)
% Compile via the wrapper plan.tex (uses tikz + positioning + shapes).

\begin{tikzpicture}[
    every node/.style={font=\small\ttfamily, align=left},
    abstract/.style={rectangle, draw, very thick, fill=blue!10,
                     minimum width=42mm, minimum height=12mm,
                     rounded corners=2pt},
    concrete/.style={rectangle, draw, fill=green!10,
                     minimum width=42mm, minimum height=12mm,
                     rounded corners=2pt},
    iface/.style={rectangle, draw, dashed, fill=yellow!15,
                  minimum width=42mm, minimum height=10mm,
                  rounded corners=2pt},
    config/.style={rectangle, draw, fill=orange!10,
                   minimum width=38mm, minimum height=10mm,
                   rounded corners=2pt},
    result/.style={rectangle, draw, fill=purple!10,
                   minimum width=38mm, minimum height=10mm,
                   rounded corners=2pt},
    inh/.style={->, >=open triangle 60, thick},   % inheritance
    impl/.style={->, >=open triangle 60, dashed, thick},   % implements
    use/.style={->, >=stealth, dotted, thick},    % uses / delegates
    comp/.style={->, >=diamond, thick},           % composition
]

% Experiment hierarchy (top)
\node[abstract] (exp) at (0, 8) {\itshape Experiment\\(ABC)};
\node[concrete] (sce) at (-5.5, 5) {SingleCell\\Experiment};
\node[concrete] (swe) at (0, 5) {Sweep\\Experiment};
\node[concrete] (kae) at (5.5, 5) {KomondorArray\\Experiment};

\draw[inh] (sce) -- (exp);
\draw[inh] (swe) -- (exp);
\draw[inh] (kae) -- (exp);

% Sweep delegates to SingleCell
\draw[use] (swe.south west) to[bend right=15]
    node[above, sloped, font=\tiny] {per-cell} (sce.east);

% Config (composition root)
\node[config] (cfg) at (0, 2.5) {ExperimentConfig};
\node[config] (cell) at (-5.5, 0.8) {CellSpec};
\node[config] (mod) at (-2.2, 0.8) {ModelConfig};
\node[config] (top) at ( 1.2, 0.8) {TopKConfig};
\node[config] (ker) at ( 4.5, 0.8) {KernelConfig};
\node[config] (cc)  at (-5.5, -0.8) {CycleCacheConfig};
\node[config] (out) at (-2.2, -0.8) {OutputConfig};

\draw[comp] (cell) -- (cfg);
\draw[comp] (mod) -- (cfg);
\draw[comp] (top) -- (cfg);
\draw[comp] (ker) -- (cfg);
\draw[comp] (cc) to[bend left=10] (cfg);
\draw[comp] (out) to[bend right=10] (cfg);

% Monitor side (right column)
\node[iface] (mon) at (10, 8) {\itshape Monitor\\(Protocol)};
\node[concrete] (wm) at (10, 6.5) {WallMonitor};
\node[concrete] (mm) at (10, 5.3) {MemoryMonitor};
\node[concrete] (am) at (10, 4.1) {AUCMonitor};
\node[concrete] (je) at (10, 2.9) {JSONLEmitter};

\draw[impl] (wm.north) -- (mon.south west);
\draw[impl] (mm.north) -- (mon.south);
\draw[impl] (am.north east) -- (mon.south east);
\draw[impl] (je.north east) to[bend right=15] (mon.east);

% Experiment uses Monitor
\draw[use, very thick] (exp.east) -- node[above, font=\tiny] {observes}
    (mon.west);

% Config -> Experiment
\draw[use] (cfg.north) -- node[right, font=\tiny] {drives} (exp.south);

% Result types (bottom right)
\node[result] (res) at (10, 1.0) {ExperimentResult};
\node[result] (sum) at (10, -0.5) {SweepSummary};

\draw[use] (sce.south east) to[bend right=10] (res.west);
\draw[use] (swe.south east) to[bend right=20] (sum.west);

% Registry (bottom)
\node[concrete] (reg) at (0, -2.5) {ExperimentRegistry\\\textit{(singleton)}};
\draw[use] (reg.north) -- node[right, font=\tiny] {YAML name $\to$ class} (cfg.south);

% Legend
\matrix[draw, fill=white, anchor=south west, font=\tiny,
        column sep=1ex, row sep=0.4ex, inner sep=4pt]
    at (-6.8, -3.5) {
    \node[abstract, minimum width=8mm, minimum height=5mm]{}; & \node{abstract / ABC}; \\
    \node[concrete, minimum width=8mm, minimum height=5mm]{}; & \node{concrete}; \\
    \node[iface, minimum width=8mm, minimum height=5mm]{}; & \node{Protocol / interface}; \\
    \node[config, minimum width=8mm, minimum height=5mm]{}; & \node{config dataclass}; \\
    \node[result, minimum width=8mm, minimum height=5mm]{}; & \node{result type}; \\
};

\matrix[draw, fill=white, anchor=south west, font=\tiny,
        column sep=1ex, row sep=0.4ex, inner sep=4pt]
    at (-1.5, -3.5) {
    \draw[inh] (0, 0) -- (0.8, 0); & \node{inheritance};   \\
    \draw[impl] (0, 0) -- (0.8, 0); & \node{implements};   \\
    \draw[use] (0, 0) -- (0.8, 0); & \node{uses / delegates}; \\
    \draw[comp] (0, 0) -- (0.8, 0); & \node{composition};  \\
};

\end{tikzpicture}

  \caption{Experiment runner framework. \textit{Italic} = abstract /
    protocol. Concrete subclasses fan out to the three deployment
    targets (single-cell, multi-seed sweep, Komondor array).
    \texttt{ExperimentConfig} is a composition of existing
    \texttt{runtime\_config.py} dataclasses
    (\texttt{TopKConfig}, \texttt{KernelConfig}, \dots), extended
    with cell / model / output specifications. Monitors are an
    Observer pattern attached to every \texttt{Experiment}; the
    registry maps YAML \texttt{experiment\_class} names to concrete
    \texttt{Experiment} subclasses.}
\end{figure}

\section*{Mermaid version}
A browsable Mermaid version of the same diagram is at
\texttt{runner\_class\_diagram.mmd} (compatible with GitHub
rendering and the user's IDE).

\end{document}
